\section{Discussion and limitations}
\label{sec:discussion}

\textit{[This section will be written after confirmatory analysis is
complete. The structure below maps to the surviving claim cards; it does
not pre-commit to interpreting results that have not been observed.]}

\subsection{Interpretation of primary results}

\textit{[Interpret H1 and H2 against their preregistered evidence
thresholds. If the warning lead-time improvement and the recoverability
decline both survive their sensitivity conditions, discuss the practical
meaning of the observed effect sizes in terms of actionable steps and
recovery opportunity. If only one passes, interpret the partial result
without claiming the other. If neither passes, discuss the benchmark value
independently of the predictive claims.]}

\subsection{Causal versus predictive language}

Results from controlled fault injection (randomized assignment with known
causal parent) support causal interpretation under the standard assumptions
of no unmeasured confounders within the randomized block. Results from
naturally occurring L4 failures are reported with associational language
only. Warning lead time is a predictive claim; the warning model observes
correlates of failure, not necessarily causes.

\subsection{Coverage limits}

All primary results are conditional on the sampled models, task families,
fault mechanisms, and resource budgets. The benchmark tests two model
configurations, four workload families, and six retained fault mechanisms
after pilot filtering. Generalization to substantially different agent
architectures, scientific domains, or fault types is not claimed without
further evidence. The external-validity L4 arm is sized after the pilot and
is explicitly a transfer test, not a population estimate.

\subsection{Falsification failures and negative findings}

All predeclared sensitivity tests that fail are reported. This includes
any failure law that does not replicate, any warning improvement that
disappears under terminal-event removal, any attribution degradation on
unseen families, and any recoverability result explained entirely by
severity rather than delay. A negative finding on a preregistered test is
retained in the artifact and discussed; it does not trigger post-hoc
outcome substitution.

\subsection{Comparison with prior benchmarks}

\textit{[Position the obtained failure-dynamics measurements against the
failure-taxonomy results in MAST-Data~\citep{mast2025}, the critical-step
localization in AgentRx~\citep{agentrx2026}, and the intervention-oriented
taxonomy in Aegis~\citep{aegis2025}. Discuss whether the temporal,
execution-grounded approach reveals structure that static or post-hoc
labeling does not.]}
